\documentclass[../../数学分析培优讲义.tex]{subfiles}

\begin{document}

\setcounter{chapter}{6}
\pagestyle{fancy}

\chapter{反常积分}


\section{积分计算}

本节基本方法前文都已介绍过,包括 N-L 公式、换元、分部积分、拟合法等等,差别仅在套上了反常积分的语言而已.



\subsection{基本积分计算}

\begin{example}
计算下列积分:
\begin{enumerate}[label=(\arabic*)]
    \item $\displaystyle\int_0^{+\infty}\dfrac{x^2}{1+x^4}\,\mathrm{d}x$;
    \item $\displaystyle\int_{-\infty}^{+\infty}|t-x|^{\frac{1}{2}}
    \dfrac{y\,\mathrm{d}t}{(t-x^2)+y^2}$;
    \item $\displaystyle\int_0^{+\infty}\dfrac{\ln x}{1+x^2}\,\mathrm{d}x$;
    \item $\displaystyle\int_0^{+\infty}\dfrac{\mathrm{d}x}{(1+x^2)(1+x^\alpha)}$;
    \item $\displaystyle\int_0^{+\infty}f(x^p+x^{-p})\dfrac{\ln x}{1+x^2}\,\mathrm{d}x$.
\end{enumerate}
\end{example}
\exampleblank

\begin{example}
计算下列积分:
\begin{enumerate}[label=(\arabic*)]
    \item $\displaystyle\int_0^{\frac{\pi}{2}}\cos(2nx)\ln(\cos x)\,\mathrm{d}x$;
    \item $\displaystyle\int_a^b\dfrac{\mathrm{d}x}{\sqrt{(x-a)(b-x)}}\quad(b>a)$;
    \item $\displaystyle\int_{-1}^{1}\dfrac{\mathrm{d}x}{(a-x)\sqrt{1-x^2}}\quad(a>1)$.
\end{enumerate}
\end{example}
\exampleblank

\begin{example}
计算下列积分:
\begin{enumerate}[label=(\arabic*)]
    \item $\displaystyle\int_{-\infty}^{+\infty}\dfrac{\mathrm{d}x}{(x^2+2x+2)^n}$;
    \item $\displaystyle\int_0^1x^n\left(\ln\dfrac{1}{x}\right)^m\,\mathrm{d}x\quad(n,m\in\mathbb{N})$.
\end{enumerate}
\end{example}
\exampleblank

\begin{example}
证明:
\begin{enumerate}[label=(\arabic*)]
    \item
    \[
        \int_0^{+\infty}f\left(ax+\dfrac{b}{x}\right)\,\mathrm{d}x
        =\dfrac{1}{a}\int_0^{+\infty}f\left(\sqrt{t^2+4ab}\right)\,\mathrm{d}t,
        \qquad a,b>0,
    \]
    假定两积分有意义;
    \item
    \[
        \int_0^{+\infty}f\left(\left(Ax-\dfrac{B}{x}\right)^2\right)\,\mathrm{d}x
        =\dfrac{1}{A}\int_0^{+\infty}f(y^2)\,\mathrm{d}y,
        \qquad A,B>0,
    \]
    假定两积分有意义.
\end{enumerate}
\end{example}
\exampleblank

\begin{example}[\markstar]
证明:
\[
    \int_1^{+\infty}\left(\dfrac{1}{[x]}-\dfrac{1}{x}\right)\,\mathrm{d}x
    =\lim\limits_{n\to\infty}\left(1+\dfrac{1}{2}+\cdots+\dfrac{1}{n}-\ln n\right)
    =\gamma,
\]
其中 $\gamma$ 为 Euler 常数.
\end{example}
\exampleblank



\subsection{特殊积分计算}

Froullani 积分、Euler 积分、Dirichlet 积分、概率积分均已在基础教材中介绍,故不再赘述,它们的证明、形式及应用均需熟练掌握.

\begin{example}
利用 Euler 积分,计算下列积分:
\begin{enumerate}[label=(\arabic*)]
    \item $\displaystyle\int_0^{\frac{\pi}{2}}\ln\tan x\,\mathrm{d}x$;
    \item $\displaystyle\int_0^1\dfrac{\ln x}{\sqrt{1-x^2}}\,\mathrm{d}x$;
    \item $\displaystyle\int_0^1\dfrac{\arcsin x}{x}\,\mathrm{d}x$;
    \item $\displaystyle\int_0^{\frac{\pi}{2}}x\cot x\,\mathrm{d}x$;
    \item $\displaystyle\int_0^\pi x\ln\sin x\,\mathrm{d}x$;
    \item $\displaystyle\int_0^\pi\dfrac{x\sin x}{1-\cos x}\,\mathrm{d}x$.
\end{enumerate}
\end{example}
\exampleblank

\begin{example}
利用 Froullani 积分,计算下列积分($a,b>0$):
\begin{enumerate}[label=(\arabic*)]
    \item $\displaystyle\int_0^{+\infty}\dfrac{\arctan ax-\arctan bx}{x}\,\mathrm{d}x$;
    \item $\displaystyle\int_0^{+\infty}\dfrac{e^{-ax}-e^{-bx}}{x}\,\mathrm{d}x$;
    \item $\displaystyle\int_0^{+\infty}\dfrac{\cos ax-\cos bx}{x}\,\mathrm{d}x$;
    \item $\displaystyle\int_0^{+\infty}\dfrac{\sin ax\sin bx}{x}\,\mathrm{d}x$;
    \item $\displaystyle\int_0^1\dfrac{x^{a-1}-x^{b-1}}{\ln x}\,\mathrm{d}x$;
    \item $\displaystyle\int_0^{+\infty}\dfrac{b\sin ax-a\sin bx}{x^2}\,\mathrm{d}x$.
\end{enumerate}
\end{example}
\exampleblank

\begin{example}
利用 Dirichlet 积分,计算下列积分:
\begin{enumerate}[label=(\arabic*)]
    \item $\displaystyle\int_0^{+\infty}\dfrac{\sin^2x}{x^2}\,\mathrm{d}x$;
    \item $\displaystyle\int_0^{+\infty}\dfrac{\sin^4x}{x^2}\,\mathrm{d}x$;
    \item $\displaystyle\int_0^{+\infty}\dfrac{\sin^4x}{x^4}\,\mathrm{d}x$;
    \item $\displaystyle\int_0^{+\infty}\dfrac{x-\sin x}{x^3}\,\mathrm{d}x$;
    \item $\displaystyle\int_{-\infty}^{+\infty}\dfrac{\sin x}{x(x-\pi)}\,\mathrm{d}x$;
    \item $\displaystyle\int_0^{+\infty}\dfrac{\sin x^2}{x}\,\mathrm{d}x$.
\end{enumerate}
\end{example}
\exampleblank

\begin{example}
利用概率积分,计算下列积分:
\begin{enumerate}[label=(\arabic*)]
    \item $\displaystyle\int_0^{+\infty}e^{-x^2}\left(x^2+\dfrac{1}{2}\right)^2\,\mathrm{d}x$;
    \item $\displaystyle\int_0^{+\infty}\left(e^{-a^2x^2}-e^{-b^2x^2}\right)\,\mathrm{d}x$.
\end{enumerate}
\end{example}
\exampleblank

\begin{example}
证明:
\begin{enumerate}[label=(\arabic*)]
    \item $1-x^2\leq e^{-x^2}\leq\dfrac{1}{1+x^2}$;
    \item 由 (1) 计算概率积分.
\end{enumerate}
\end{example}
\exampleblank



\subsection{积分与极限}

该类问题大多利用拟合法结合反常积分收敛的定义或性质分段分析即可.

\begin{example}
设 $f:\mathbb{R}\to\mathbb{R}$ 是有界连续函数,求
\[
    \lim\limits_{t\to0^+}\int_{-\infty}^{+\infty}\dfrac{t}{t^2+x^2}f(x)\,\mathrm{d}x.
\]
\end{example}
\exampleblank

\begin{example}
求极限
\[
    \lim\limits_{x\to0^+}x^\alpha\int_x^1\dfrac{f(t)}{t^{\alpha+1}}\,\mathrm{d}t,
\]
其中 $\alpha>0$,$f(x)\in C\lbrack 0,1\rbrack$.
\end{example}
\exampleblank

\begin{example}
设当 $x>a$ 时,$\varphi(x)>0$,$f(x),\varphi(x)$ 在任何有限区间 $\lbrack a,b\rbrack$ 上可积,
\[
    \int_a^{+\infty}\varphi(t)\,\mathrm{d}t
\]
发散,当 $x\to+\infty$ 时,$f(x)=o(\varphi(x))$.证明:
\[
    \int_a^{+\infty}f(t)\,\mathrm{d}t
    =o\left(\int_a^{+\infty}\varphi(t)\,\mathrm{d}t\right).
\]
\end{example}
\exampleblank

\begin{proposition}[\markstar\ Riemann 引理]
设 $f(x)$ 在 $\lbrack a,+\infty)$ 上绝对可积,$g(x)$ 是周期为 $T$ 的函数,在 $\lbrack 0,T\rbrack$ 上正常可积,则
\[
    \lim\limits_{n\to\infty}\int_a^{+\infty}f(x)g(nx)\,\mathrm{d}x
    =\dfrac{1}{T}\int_0^Tg(x)\,\mathrm{d}x\int_a^{+\infty}f(x)\,\mathrm{d}x.
\]
\end{proposition}

\begin{example}
\begin{enumerate}[label=(\arabic*)]
    \item 设 $\varphi(x)$ 为有界周期函数,周期为 $T$,且
    \[
        \dfrac{1}{T}\int_0^T\varphi(x)\,\mathrm{d}x=C,
    \]
    求证:
    \[
        \lim\limits_{n\to\infty}n\int_0^{+\infty}\dfrac{\varphi(t)}{t^2}\,\mathrm{d}t=C;
    \]
    \item 设 $\displaystyle\varphi(x)=\int_0^x\cos\dfrac{1}{t}\,\mathrm{d}t$,求 $\varphi'(0)$;
    \item 设 $\displaystyle F(x)=\int_0^x\left(\dfrac{1}{t}-\left[\dfrac{1}{t}\right]\right)\,\mathrm{d}t$,试证 $F'(0)=\dfrac{1}{2}$.
\end{enumerate}
\end{example}
\exampleblank

\begin{example}
设 $f(x)$ 在任一有限区间 $\lbrack 0,a\rbrack$($a>0$)上正常可积,在 $(0,+\infty)$ 上绝对可积,证明:
\[
    \lim\limits_{n\to\infty}\int_0^{+\infty}f(x)|\sin nx|\,\mathrm{d}x
    =\dfrac{2}{\pi}\int_0^{+\infty}f(x)\,\mathrm{d}x.
\]
\end{example}
\exampleblank

\begin{example}
设 $p(t)$ 为 $\mathbb{R}$ 上连续函数满足:
\begin{enumerate}[label=(\arabic*)]
    \item $p(t)=0$,$|t|\geq1$;
    \item $\displaystyle\int_{-\infty}^{+\infty}p(t)\,\mathrm{d}t=0$;
    \item $\displaystyle\int_{-\infty}^{+\infty}tp(t)\,\mathrm{d}t=1$.
\end{enumerate}
又 $f(x)$ 为 $\mathbb{R}$ 上可微函数,证明:
\[
    \lim\limits_{\lambda\to0}\int_{-\infty}^{+\infty}
    \dfrac{1}{\lambda^2}p\left(\dfrac{t-x}{\lambda}\right)f(t)\,\mathrm{d}t
    =f'(x).
\]
\end{example}
\exampleblank

我们知道,反常积分不能像常义积分那样定义为积分和的极限.接下来这个命题说明:对于单调无界反常积分,我们可以用类似于积分和的和式来逼近.

\begin{proposition}[\markstar]
设 $f(x)$ 在 $(0,1)$ 上单调,$x=0,1$ 为奇点,$\displaystyle\int_0^1f(x)\,\mathrm{d}x$ 存在.证明:
\[
    \lim\limits_{n\to\infty}
    \dfrac{f\left(\dfrac{1}{n}\right)+f\left(\dfrac{2}{n}\right)+\cdots+f\left(\dfrac{n-1}{n}\right)}{n}
    =\int_0^1f(x)\,\mathrm{d}x.
\]
\end{proposition}

\begin{remark}
\begin{enumerate}[label=\textcircled{\arabic*}]
    \item 单调性条件可减弱为只在奇点某邻域内单调;
    \item 该结论还可推广为更一般的情况,$f$ 条件同上,$\varphi(x)$ 在 $\lbrack 0,1\rbrack$ 上常义可积,则
    \[
        \lim\limits_{n\to\infty}\dfrac{1}{n}\sum\limits_{i=1}^{n-1}
        \varphi\left(\dfrac{i}{n}\right)f\left(\dfrac{i}{n}\right)
        =\int_0^1\varphi(x)f(x)\,\mathrm{d}x;
    \]
    \item 逆命题不成立,例如 $f(x)=\dfrac{1}{x}-\dfrac{1}{1-x}$.
\end{enumerate}
\end{remark}

\begin{example}
设 $0<a<d$,
\[
    A_n=\dfrac{a+(a+d)+\cdots+[a+(n-1)d]}{n},
\]
\[
    G_n=\sqrt[n]{a(a+d)\cdots[a+(n-1)d]}.
\]
试证:
\[
    \lim\limits_{n\to\infty}\dfrac{G_n}{A_n}=\dfrac{2}{e}.
\]
\end{example}
\exampleblank

\begin{example}
$\{C_{n}^{k}\}_{k=0}^n$ 为二项式系数,$A_n,G_n$ 分别表示它们的算术平均值与几何平均值,试证:
\[
    \lim\limits_{n\to\infty}\sqrt[n]{A_n}=2,
    \qquad
    \lim\limits_{n\to\infty}\sqrt[n]{G_n}=\sqrt{e}.
\]
\end{example}
\exampleblank

\begin{example}
设 $f(x)$ 在 $\lbrack 1,+\infty)$ 上可微,$f(1)=1$,若有
\[
    f'(x)=\dfrac{1}{x^2+f^2(x)},
    \qquad 1\leq x<\infty,
\]
则
\[
    \lim\limits_{x\to+\infty}f(x)<1+\dfrac{\pi}{4}.
\]
\end{example}
\exampleblank


\newpage

\section{敛散性判别}

\begin{remark}
判别反常积分敛散性的常用方法:
\begin{enumerate}[label=(\arabic*)]
    \item 设当 $x$ 充分大时 $f(x)\geq0$:
    \begin{enumerate}[label=\textcircled{\arabic*}]
        \item 若 $\displaystyle\lim\limits_{x\to+\infty}f(x)=0$,可考察 $x\to+\infty$ 时无穷小量 $f(x)$ 的阶:若阶数 $\lambda>1$,则收敛;当 $\lambda\leq1$ 时,发散;
        \item 利用比较判别法或其极限形式;
        \item 可考察 $\displaystyle\int_a^Af(x)\,\mathrm{d}x$ 是否有界.
    \end{enumerate}
    \item 设 $x\to+\infty$ 时,$f(x)$ 无穷次变号,则考虑用 Abel 或 Dirichlet 判别法;
    \item 若前述方法均失效,可考虑用 Cauchy 收敛准则或定义;
    \item 利用分部积分或换元转化为熟悉或容易判断的积分;
    \item 利用级数(下一章详细介绍).
\end{enumerate}
\end{remark}



\subsection{非负函数积分敛散性判别}

\begin{example}
讨论下列反常积分的敛散性:
\begin{enumerate}[label=(\arabic*)]
    \item $\displaystyle\int_0^{+\infty}\dfrac{\mathrm{d}x}{\sqrt[3]{x^4+1}}$;
    \item $\displaystyle\int_0^{\frac{\pi}{2}}\dfrac{1-\cos x}{x^p}\,\mathrm{d}x$;
    \item $\displaystyle\int_0^{+\infty}\dfrac{x\arctan x}{1+x^p}\,\mathrm{d}x$;
    \item $\displaystyle\int_0^{+\infty}\dfrac{\mathrm{d}x}{x^p+x^q}$;
    \item $\displaystyle\int_0^{+\infty}\dfrac{x^p}{1+x^q}\,\mathrm{d}x$.
\end{enumerate}
\end{example}
\exampleblank

\begin{example}
讨论下列积分的敛散性:
\begin{enumerate}[label=(\arabic*)]
    \item $\displaystyle\int_1^{+\infty}\dfrac{\mathrm{d}x}{\sqrt[r]{\ln x}}\quad(r>0)$;
    \item $\displaystyle\int_2^{+\infty}\dfrac{\mathrm{d}x}{(\ln x)^{\ln x}}$;
    \item $\displaystyle\int_3^{+\infty}\dfrac{\mathrm{d}x}{(\ln x)^{\ln\ln x}}$;
    \item $\displaystyle\int_2^{+\infty}\dfrac{\mathrm{d}x}{(\ln x)(\ln x)^\alpha}$.
\end{enumerate}
\end{example}
\exampleblank

\begin{example}
如果反常积分 $\displaystyle\int_a^bf^2$ 收敛,则称 $f$ 在 $\lbrack a,b\rbrack$ 上平方可积.分无穷积分与瑕积分两种情况讨论平方可积性与绝对可积性之间的关系.
\end{example}
\exampleblank

\begin{example}
讨论
\[
    \int_0^\pi\dfrac{\sin^{\alpha-1}x}{|1+k\cos x|^\alpha}\,\mathrm{d}x
\]
的敛散性.
\end{example}
\exampleblank

\begin{example}
设 $f(x)$ 在 $\lbrack 1,+\infty)$ 上连续,对任意 $x\in\lbrack 1,+\infty)$ 有 $f(x)>0$,且
\[
    \lim\limits_{x\to+\infty}\dfrac{\ln f(x)}{\ln x}=\lambda.
\]
证明:$\lambda>1$ 时,$\displaystyle\int_1^{+\infty}f(x)\,\mathrm{d}x$ 收敛.
\end{example}
\exampleblank

\begin{example}
设 $f(x)$ 在 $(-\infty,+\infty)$ 上有定义,$f(x)>0$,且在任意有限区间 $\lbrack -A,B\rbrack$ 上可积($A,B>0$);又有定数 $M$,使得
\[
    \int_{-\infty}^{+\infty}f(x)e^{-\frac{|x|}{k}}\,\mathrm{d}x<M
\]
对任意 $k>0$ 成立.证明:$\displaystyle\int_{-\infty}^{+\infty}f(x)\,\mathrm{d}x$ 收敛.
\end{example}
\exampleblank

\begin{example}
设函数 $f(x)$ 在 $\lbrack 0,1\rbrack$ 上连续,且 $\displaystyle\lim\limits_{x\to0}f(x)=\pm\infty$.对任意 $N\in\mathbb{Z}^+$,令
\[
    f_N=\min\{f(x),N\}.
\]
证明:
\[
    \int_0^1f(x)\,\mathrm{d}x\text{ 收敛}
    \quad\Longleftrightarrow\quad
    \lim\limits_{N\to+\infty}\int_0^1f_N(x)\,\mathrm{d}x\text{ 存在}.
\]
\end{example}
\exampleblank

\begin{example}
设 $f(x)$ 是 $\lbrack 1,+\infty)$ 上的正值函数,若 $\displaystyle\int_1^{+\infty}f(x)\,\mathrm{d}x$ 收敛,则:
\begin{enumerate}[label=(\arabic*)]
    \item $\displaystyle\int_1^{+\infty}\dfrac{\sqrt{f(x)}}{x^p}\,\mathrm{d}x$($p>\dfrac{1}{2}$)收敛;
    \item $\displaystyle\int_1^{+\infty}\dfrac{\mathrm{d}x}{x^2f(x)}$ 发散;
    \item $\displaystyle\int_1^{+\infty}\sqrt{f(x)f(x+x_0)}\,\mathrm{d}x$ 收敛.
\end{enumerate}
\end{example}
\exampleblank

\begin{example}
\begin{enumerate}[label=(\arabic*)]
    \item 设 $f\in C\lbrack 1,+\infty)$,$f(x)>0$,$\displaystyle F(x)=\int_1^xf(t)\,\mathrm{d}t$,若 $\displaystyle\int_1^{+\infty}f(x)\,\mathrm{d}x=+\infty$,则
    \[
        \int_1^{+\infty}\dfrac{f(x)}{F(x)\ln F(x)}\,\mathrm{d}x
    \]
    发散,而
    \[
        \int_2^{+\infty}\dfrac{f(x)}{F(x)\ln^2F(x)}\,\mathrm{d}x
    \]
    收敛;
    \item 设 $F(x)$ 在 $\lbrack 1,+\infty)$ 上递减趋于 $0$,则 $\displaystyle\int_1^{+\infty}f(x)\,\mathrm{d}x$ 与 $\displaystyle\int_1^{+\infty}a^xf(a^x)\,\mathrm{d}x$($a>1$)同敛散.
\end{enumerate}
\end{example}
\exampleblank

\begin{example}
设 $f(x)$ 在 $\lbrack 0,+\infty)$ 上是正值递增函数,$\displaystyle F(x)=\int_0^xf(t)\,\mathrm{d}t$,$\displaystyle\lim\limits_{x\to0^+}\dfrac{x}{F(x)}=1$,则
\[
    \int_0^{+\infty}\dfrac{1}{f(x)}\,\mathrm{d}x\text{ 收敛}
    \quad\Longleftrightarrow\quad
    \int_0^{+\infty}\dfrac{x}{F(x)}\,\mathrm{d}x\text{ 收敛}.
\]
\end{example}
\exampleblank

\begin{example}
设 $f\in C^{(2)}\lbrack a,+\infty)$,若积分
\[
    \int_a^{+\infty}f^2(x)\,\mathrm{d}x,
    \qquad
    \int_a^{+\infty}[f''(x)]^2\,\mathrm{d}x
\]
均收敛,则 $\displaystyle\int_a^{+\infty}[f'(x)]^2\,\mathrm{d}x$ 收敛.
\end{example}
\exampleblank

\begin{example}
设 $f(x)>0$,$\displaystyle\int_1^{+\infty}f(x)\,\mathrm{d}x$ 收敛.证明:$[f(x)]^{1-\frac{1}{x}}$ 收敛.
\end{example}
\exampleblank

\begin{example}
\begin{enumerate}[label=(\arabic*)]
    \item 设 $f(x)>0$($a\leq x<+\infty$),且对任意 $b>a$ 有 $f\in R\lbrack a,b\rbrack$.若对 $x_0>a$,有
    \[
        \lim\limits_{x\to+\infty}\dfrac{f(x+x_0)}{f(x)}=l,
    \]
    则
    \[
        \int_a^{+\infty}f(x)\,\mathrm{d}x
        \begin{cases}
            \text{收敛},&l<1,\\
            \text{发散},&l>1;
        \end{cases}
    \]
    \item 设 $f\in R\lbrack a,b\rbrack$($\forall b>a$),且 $f(x)>0$($a\leq x<+\infty$),
    \[
        \lim\limits_{x\to+\infty}\dfrac{f(kx)}{f(x)}=l,
        \qquad k>1,
    \]
    则
    \[
        \int_a^{+\infty}f(x)\,\mathrm{d}x
        \begin{cases}
            \text{收敛},&l<\dfrac{1}{k},\\
            \text{发散},&l>\dfrac{1}{k}.
        \end{cases}
    \]
\end{enumerate}
\end{example}
\exampleblank

类似于级数中判别法的讨论,比较判别法中常用的比较对象是 $p$ 积分 $\displaystyle\int_1^{+\infty}\dfrac{\mathrm{d}x}{x^p}$.那么,是否有一个普遍适用的比较标准呢?回答是否定的.

\begin{proposition}[\markstar]
设 $f(x)\in C\lbrack a,+\infty)$ 且 $f(x)>0$,则:
\begin{enumerate}[label=(\arabic*)]
    \item 若 $\displaystyle\int_a^{+\infty}f(x)\,\mathrm{d}x$ 发散,令 $\displaystyle F(x)=\int_a^xf(t)\,\mathrm{d}t$,则
    \begin{enumerate}[label=\textcircled{\arabic*}]
        \item $\displaystyle\int_a^{+\infty}\dfrac{f(x)}{F(x)}\,\mathrm{d}x$ 发散;
        \item $\displaystyle\int_a^{+\infty}\dfrac{f(x)}{F^{1+\delta}(x)}\,\mathrm{d}x$($\delta>0$)收敛.
    \end{enumerate}
    \item 若 $\displaystyle\int_a^{+\infty}f(x)\,\mathrm{d}x$ 收敛,令 $\displaystyle R(x)=\int_x^{+\infty}f(t)\,\mathrm{d}t$,则
    \begin{enumerate}[label=\textcircled{\arabic*}]
        \item $\displaystyle\int_a^{+\infty}\dfrac{f(x)}{R(x)}\,\mathrm{d}x$ 发散;
        \item $\displaystyle\int_a^{+\infty}\dfrac{f(x)}{R^{1-\delta}(x)}\,\mathrm{d}x$($\delta>0$)收敛.
    \end{enumerate}
\end{enumerate}
\end{proposition}



\subsection{一般函数积分敛散性判别}

\begin{proposition}[\markstar]
\begin{enumerate}[label=(\arabic*)]
    \item $\displaystyle\int_1^{+\infty}\dfrac{\sin x}{x^p}\,\mathrm{d}x$ 与 $\displaystyle\int_1^{+\infty}\dfrac{\cos x}{x^p}\,\mathrm{d}x$ 敛散性相同,且
    \[
        \begin{cases}
            \text{绝对收敛},&p>1,\\
            \text{条件收敛},&0<p\leq1,\\
            \text{发散},&p\leq0;
        \end{cases}
    \]
    \item
    \[
        \int_0^\infty\dfrac{\sin x}{x^p}\,\mathrm{d}x
        \begin{cases}
            \text{绝对收敛},&1<p<2,\\
            \text{条件收敛},&0<p\leq1,\\
            \text{发散},&\text{其它},
        \end{cases}
    \]
    而
    \[
        \int_0^{+\infty}\dfrac{\cos x}{x^p}\,\mathrm{d}x
        \begin{cases}
            \text{条件收敛},&0<p<1,\\
            \text{发散},&\text{其它}.
        \end{cases}
    \]
\end{enumerate}
\end{proposition}

\begin{example}
证明:积分
\[
    \int_0^1\left(x\sin\dfrac{1}{x^2}-\dfrac{1}{x}\cos\dfrac{1}{x^2}\right)\,\mathrm{d}x
\]
收敛.
\end{example}
\exampleblank

\begin{example}
讨论下列反常积分的条件收敛与绝对收敛性:
\begin{enumerate}[label=(\arabic*)]
    \item $\displaystyle\int_0^{+\infty}\sin x^2\,\mathrm{d}x$;
    \item $\displaystyle\int_0^{+\infty}\cos^2x\,\mathrm{d}x$;
    \item $\displaystyle\int_0^{+\infty}x\sin^4x\,\mathrm{d}x$;
    \item $\displaystyle\int_0^1\dfrac{1}{x^p}\sin\dfrac{1}{x}\,\mathrm{d}x$.
\end{enumerate}
\end{example}
\exampleblank

\begin{example}
积分
\[
    \int_0^{+\infty}\left[\left(1-\dfrac{\sin x}{x}\right)^{-\frac{1}{3}}-1\right]\,\mathrm{d}x
\]
是否收敛?是否绝对收敛?
\end{example}
\exampleblank

\begin{example}
设 $\alpha,\beta\in\mathbb{R}$,讨论积分
\[
    \int_0^{+\infty}x^\alpha\sin(x^\beta)\,\mathrm{d}x
\]
的绝对收敛性与条件收敛性.
\end{example}
\exampleblank

\begin{example}
讨论下列积分的绝对收敛性与条件收敛性:
\begin{enumerate}[label=(\arabic*)]
    \item $\displaystyle\int_0^{+\infty}\dfrac{\sin\left(x+\dfrac{1}{x}\right)}{x^\alpha}\,\mathrm{d}x$;
    \item $\displaystyle\int_1^{+\infty}\sin\left(\dfrac{\sin x}{\sqrt{x}}\right)\dfrac{\mathrm{d}x}{\sqrt{x}}$.
\end{enumerate}
\end{example}
\exampleblank

\begin{example}
判断下列积分的敛散性:
\begin{enumerate}[label=(\arabic*)]
    \item $\displaystyle\int_1^{+\infty}\tan\left(\dfrac{\sin x}{x}\right)\,\mathrm{d}x$;
    \item $\displaystyle\int_1^{+\infty}\dfrac{\ln\cos\dfrac{1}{x}}{x^p}\,\mathrm{d}x$;
    \item $\displaystyle\int_1^{+\infty}\dfrac{\sin x}{x^\alpha-\sin x}\,\mathrm{d}x$.
\end{enumerate}
\end{example}
\exampleblank

\begin{example}
设 $f(x)>0$ 单调递减,试证:
\[
    \int_a^{+\infty}f(x)\,\mathrm{d}x
    \quad\text{与}\quad
    \int_a^{+\infty}f(x)\sin^2x\,\mathrm{d}x
\]
同敛散.
\end{example}
\exampleblank

\begin{example}
设 $f(x)$ 在 $\lbrack a,+\infty)$ 上可微,且 $x\to+\infty$ 时 $f'(x)$ 单调递增趋于 $+\infty$.证明:
\[
    \int_a^{+\infty}\sin(f(x))\,\mathrm{d}x
    \quad\text{和}\quad
    \int_a^{+\infty}\cos(f(x))\,\mathrm{d}x
\]
都收敛.
\end{example}
\exampleblank

对于周期函数特有的积分性质 $\displaystyle\int_a^{a+T}f(x)\,\mathrm{d}x=\int_0^Tf(x)\,\mathrm{d}x$($T$ 为 $f$ 周期),我们常对周期函数反常积分作周期处理.

\begin{example}
设 $f\in C\lbrack a,+\infty)$ 且是周期为 $T$ 的周期函数,又 $g(x)$ 是 $\lbrack a,+\infty)$ 上的单调函数,且 $\displaystyle\lim\limits_{x\to+\infty}g(x)=0$,则:
\begin{enumerate}[label=(\arabic*)]
    \item 若 $\displaystyle\int_a^{a+T}f(x)\,\mathrm{d}x=0$,则 $\displaystyle\int_a^{+\infty}f(x)g(x)\,\mathrm{d}x$ 收敛;
    \item 若 $\displaystyle\int_a^{a+T}f(x)\,\mathrm{d}x\neq0$,则 $\displaystyle\int_a^{+\infty}f(x)g(x)\,\mathrm{d}x$ 与 $\displaystyle\int_a^{+\infty}g(x)\,\mathrm{d}x$ 同敛散.
\end{enumerate}
\end{example}
\exampleblank

\begin{example}[\markstar]
讨论积分
\[
    \int_0^{+\infty}\dfrac{x}{1+x^6\sin^2x}\,\mathrm{d}x
\]
的敛散性.
\end{example}
\exampleblank

\begin{example}[\markstar]
讨论积分
\[
    \int_0^{+\infty}\dfrac{x}{1+x^4\sin^2x}\,\mathrm{d}x
\]
的敛散性.
\end{example}
\exampleblank


\newpage

\section{无穷远处性质}

反常积分许多性质可类比于级数,有时甚至完全一样,但有时也天差地别.考虑收敛级数 $\displaystyle\sum\limits_{n=1}^{\infty}a_n$,有 $\displaystyle\lim\limits_{n\to\infty}a_n=0$.自然的一个疑问是,能否有如下结论:
\[
    \int_a^{+\infty}f(x)\,\mathrm{d}x\text{ 收敛}
    \quad\Longrightarrow\quad
    \lim\limits_{x\to+\infty}f(x)=0?
\]
答案是否定的,上节最后几个例子均为反例,并且一般情况下都得不到如上结论.下面给出一个特殊情况下的成立条件.

\begin{proposition}[\markstar]
设 $\displaystyle\int_a^{+\infty}f(x)\,\mathrm{d}x$ 收敛且 $f\in C\lbrack a,+\infty)$,则
\[
    \lim\limits_{x\to+\infty}f(x)=0
    \quad\Longleftrightarrow\quad
    f\text{ 在 }\lbrack a,+\infty)\text{ 上一致连续}.
\]
\end{proposition}

\begin{corollary}[\markstar]
设 $\displaystyle\int_a^{+\infty}f(x)\,\mathrm{d}x$ 收敛,且 $|f'(x)|\leq M$,则
\[
    \lim\limits_{x\to+\infty}f(x)=0.
\]
\end{corollary}

\begin{example}
设 $\displaystyle\int_a^{+\infty}f(x)\,\mathrm{d}x$ 收敛,且 $\displaystyle\lim\limits_{x\to+\infty}f(x)$ 存在,则
\[
    \lim\limits_{x\to+\infty}f(x)=0.
\]
\end{example}
\exampleblank

\begin{example}
\begin{enumerate}[label=(\arabic*)]
    \item 设 $\displaystyle\int_a^{+\infty}f(x)\,\mathrm{d}x$ 收敛,且 $f$ 单调,则 $\displaystyle\lim\limits_{x\to+\infty}xf(x)=0$.逆命题是否成立?
    \item 设 $\displaystyle\int_a^{+\infty}x^\alpha f(x)\,\mathrm{d}x$ 收敛,且 $f$ 单调,则 $\displaystyle\lim\limits_{x\to+\infty}x^{\alpha+1}f(x)=0$.
    \item 设 $\displaystyle\int_a^{+\infty}f(x)\,\mathrm{d}x$ 收敛,且 $xf(x)$ 单调,则 $\displaystyle\lim\limits_{x\to+\infty}xf(x)\ln x=0$.
    \item 设 $\displaystyle\int_a^x f(x)\,\mathrm{d}x$ 收敛,且 $\dfrac{f(x)}{x}$ 单调递减($a>0$),则 $\displaystyle\lim\limits_{x\to+\infty}xf(x)=0$.
\end{enumerate}
\end{example}
\exampleblank

\begin{example}
\begin{enumerate}[label=(\arabic*)]
    \item 设 $f\in C\lbrack a,+\infty)$ 且 $\displaystyle\int_a^{+\infty}f(x)\,\mathrm{d}x$ 收敛,则存在数列 $\{x_n\}\subseteq\lbrack a,+\infty)$ 满足
    \[
        \lim\limits_{n\to\infty}x_n=+\infty,
        \qquad
        \lim\limits_{n\to\infty}f(x_n)=0;
    \]
    \item 设 $f\in C\lbrack a,+\infty)$ 且 $\displaystyle\int_a^{+\infty}f(x)\,\mathrm{d}x$ 收敛,则存在数列 $\{x_n\}\subseteq\lbrack a,+\infty)$ 满足
    \[
        \lim\limits_{n\to\infty}x_n=+\infty,
        \qquad
        \lim\limits_{n\to\infty}x_nf(x_n)=0;
    \]
    \item 设 $f$ 在 $\lbrack a,+\infty)$ 上可微且 $\displaystyle\int_a^{+\infty}f(x)\,\mathrm{d}x$ 收敛,证明:存在数列 $\{x_n\}\subseteq\lbrack a,+\infty)$ 满足
    \[
        \lim\limits_{n\to\infty}x_n=+\infty,
        \qquad
        \lim\limits_{n\to\infty}f'(x_n)=0.
    \]
\end{enumerate}
\end{example}
\exampleblank

\begin{example}
证明:若 $f$ 连续可微,$\displaystyle\int_a^{+\infty}f(x)\,\mathrm{d}x$ 与 $\displaystyle\int_a^{+\infty}f'(x)\,\mathrm{d}x$ 都收敛,则
\[
    \lim\limits_{x\to+\infty}f(x)=0.
\]
\end{example}
\exampleblank

\end{document}
